\documentclass[12pt,nonacm]{article}
%
\usepackage{amsmath,amssymb}
\usepackage[dvipsnames]{xcolor}
\usepackage{mathpazo}
\title{Oligo-poly\\
\large
A library for working with polynomial rings under oligomorphic group actions.}
%
\author{Arka Ghosh}
%
\newcommand{\struct}{\mathcal{X}}
\newcommand{\dloQ}{\mathcal{Q}}

\newcommand{\aut}[1]{\mathsf{Aut}\left(#1\right\right)}
\newcommand{\poly}[2]{#1[#2]}

\newcommand{\emphdef}[1]{\textcolor{BrickRed}{#1}}
%
\newtheorem{theorem}{Theorem}
\newtheorem{example}[theorem]{Example}
%
\begin{document}
\maketitle
%
The most recent version of this document can be found at...

The purpose of this library is to implement (in Rust) the equivariant Buchberger algorithm given in \cite{GL25}.

We assume that the reader is .....
%
\paragraph*{Structure of this document}

%
\section{Theoretical background}
%

\subsection{Relational structures}
%
A \emphdef{relational structure} $\struct = (X,R_1,\dots,R_n)$ consists of a
set $X$ (called the \emphdef{domain} of the relational structure) and relations $R_i$ defined on the set $X$.
%
\begin{example}
Let $Q$ be the set of rational numbers.
We can consider the structure $\dloQ = (Q,<)$ where is $<$ the usual ordering.
\end{example}
%
\begin{example}
A graph $\mathcal{G} = (V,E)$ is another example of a relational structure,
with the set $V$ of vertices being the domain and the edge relation $E$ being the only relation.
\end{example}

By abuse of notation we sometime use $\mathcal{X}$ to denote the set underlying $X$.
For instance,
we write $\bar{x}\in\struct^n$ to mean that $x$ is an element of $X^n$.

\subsection{Automorphism groups}

An \emphdef{automorphism} of the relational structure $\struct$ is a bijection from $\struct$ to itself that preserves and reflects the structure.
%
\begin{example}
An automorphism of $\dloQ$ is an order preserving bijection from $\mathcal{Q}$ to itself.
\end{example}
%
\begin{example}
If $\struct$ is a graph,
then an automorphism of $\mathcal{X}$ is an isomorphism of the graph with itself.
\end{example}
%
The set of automorphisms of a relational structure $\mathcal{X}$ is denoted as \emphdef{$\mathbf{Aut}(\mathcal{X})$}.
%
\subsection{Equivariant Ideals}
%
For a relational structure $\mathcal{X}$,
we use $\mathbb{Q}[\mathcal{X}]$ to denote the polynomial ring whose variables are elements of $\mathcal{X}$.

Instead of using $\mathcal{X}$ as the set of variables some articles use variables $a_{x}$ indexed by elements $x\in\mathcal{X}$.
But using $\mathcal{X}$ itself makes the notation simpler.

The group $\mathbf{Aut}(\mathcal{X})$ of automorphisms of $\mathcal{X}$ acts on $\mathbb{Q}[\mathcal{X}]$ in a variable-wise manner.
For instance, for $\pi\in\mathbf{Aut}(\mathcal{X})$ and $x,y\in\mathcal{X}$
$$\pi(x^2 - 2xy^2 + 3) =
\pi(x)^2 - 2\pi(x)\pi(y)^2 + 3\ .$$

An ideal in 

\subsection{Equivariant Gröbner bases}

\subsection{Equivariant Buchberger algorithm}

\section{How to use}

\subsection{Relational structures}

\subsection{Pre-defined relational structures}

\paragraph{Dense linear order without endpoints}

\subsection{Defining new relational structures}

\begin{thebibliography}{9}
\bibitem{GL25}
Arka Ghosh and Aliaume Lopez. Computability of equivariant
Gr\"{o}bner bases, 2025.

\end{thebibliography}
%
\end{document}
